% docs/style/kopf-fusszeilen.tex -- local copy of tex-skripten/formatierungen/kopf-fusszeilen.tex
% Copied 2026-10-12 so that the repository compiles stand-alone (GitHub).
% Local version: 1.0.0 -- running head shows the project name, see marker "interferosim"

% 1. Alle vorherigen Stile löschen
\clearpairofpagestyles

% 2. Linien konfigurieren (modernere Syntax für scrlayer-scrpage)
\KOMAoptions{
	headsepline=0.4pt,
	footsepline=0.4pt,
	plainheadsepline=true, % Linie auch auf Kapitalseiten oben
	plainfootsepline=true  % Linie auch auf Kapitalseiten unten
}

% 3. Inhalte definieren
% Das optionale Argument [..] definiert den Stil für 'plain' (Kapitelseiten)
% Das normale Argument {..} definiert den Stil für 'scrheadings' (Folgeseiten)

\ihead[\headmark]{\headmark}       % Innen: Kapitelname
%\ohead[\pagemark]{TAP-1}
\ohead[\pagemark]{Interferosim}   % interferosim: project name       % Außen: Seitenzahl oben (falls gewünscht)

\ifoot[Erich]{E.Schubert}               % Fuß links: Dein Name
\ofoot[\pagemark]{\pagemark}       % Fuß rechts: Seitenzahl

% 4. Den Stil global aktivieren
\pagestyle{scrheadings}

% 5. Erzwingen, dass 'plain' (Kapitelseiten) mit Linien und Text gefüllt wird
% Dies überschreibt das Standard-Verhalten von KOMA
\renewcommand*{\chapterpagestyle}{scrheadings}